\documentclass[12pt,a4paper]{article}

\usepackage[a4paper,left=2.0cm,right=2.0cm,top=2.8cm,bottom=2.5cm]{geometry}
\usepackage{fontspec}
\usepackage{enumitem}
\usepackage{microtype}
\usepackage{setspace}
\usepackage{titlesec}

\setmainfont{Times New Roman}
\setstretch{1.3}
\setlength{\parindent}{0pt}
\setlength{\parskip}{0.55em}
\setlist[itemize]{leftmargin=2em,itemsep=0.35em,topsep=0.35em}
\setlist[itemize,2]{label=--,leftmargin=2.2em}

\titleformat{\section}{\Large\bfseries}{\thesection}{1em}{}
\titleformat{\subsection}{\large\bfseries}{\thesubsection}{1em}{}
\titlespacing*{\section}{0pt}{1.1em}{0.45em}
\titlespacing*{\subsection}{0pt}{0.9em}{0.35em}

\begin{document}

\begin{center}
  {\LARGE Machine Learning Course 2026 Autumn: Experiment 1}

  \vspace{1.35cm}
  {\large September 29, 2026}
\end{center}

\vspace{0.7cm}

\section{Overview}

In this experiment, you will implement linear models, including linear regression,
ridge regression, and logistic regression, and evaluate trained models on the test
set. You will also compare your implementation with the models in
\textbf{\textit{sklearn}}.

\section{Dataset Description}

\begin{sloppypar}
The datasets are ready for you and are stored in the files
\textbf{\textit{diabetes\_train.csv}},
\textbf{\textit{diabetes\_test.csv}},
\textbf{\textit{pendigit\_train.csv}}, and
\textbf{\textit{pendigit\_test.csv}}. Each row in the files is a sample, with the
last column being the label and the other columns being the features.
\end{sloppypar}

\section{Task Description}

\subsection{Task 1}

Given the \textbf{\textit{diabetes}} dataset, implement the
\textbf{linear regression} models for \textbf{gradient descent} and
\textbf{closed-form solution}, and evaluate the trained model on the test set.

\subsection{Task 2}

Given the \textbf{\textit{diabetes}} dataset, implement the
\textbf{ridge regression} model with \textbf{gradient descent}, and evaluate the
trained model on the test set. \textbf{5-fold cross-validation} will be employed to
select the proper value of the hyperparameter, and the influence of these values on
the trained model will be evaluated. In addition, use the
\textbf{closed-form solution} to construct the \textbf{ridge regression} model using
the \textbf{optimal hyperparameter} obtained by cross-validation.

\subsection{Task 3}

Given the \textbf{\textit{pendigit}} dataset, implement the
\textbf{logistic regression} model with \textbf{gradient descent}, and evaluate the
trained model on the test set.

\subsection{Task 4}

Compare the trained models with standard models in \textbf{\textit{sklearn}}.
Discuss the difference between them and try to analyze the possible reasons.

\section{Requirements}

\begin{itemize}
  \item Implement learning algorithms with \textbf{Python}. Basic data analysis
        libraries such as \textbf{NumPy}, \textbf{Pandas}, \textbf{Matplotlib},
        etc., are allowed to use.

  \item For Tasks 1--3, you cannot use highly integrated libraries such as
        \textbf{sklearn}, \textbf{PyTorch}, \textbf{TensorFlow}, etc. For Task 4
        only, you may use \textbf{sklearn} for comparison.

  \item Record the procedure of experiments in \textbf{Jupyter Notebook}. It is
        recommended to use \textbf{Markdown} in Notebook to properly explain the
        procedure.

  \item Your submission includes:
        \begin{itemize}
          \item source code: python notebook (\textit{.ipynb});
          \item \textit{report.doc} file: your report. Refer to the report template
                for the format.
        \end{itemize}
\end{itemize}

\end{document}
